\documentclass[11pt]{article}
\usepackage[margin=1in]{geometry}
\usepackage{amsmath,amssymb}
\usepackage{booktabs}
\usepackage{enumitem}
\usepackage{hyperref}
\usepackage{xcolor}
\usepackage{listings}

\lstset{
  basicstyle=\ttfamily\footnotesize,
  breaklines=true,
  frame=single,
  numbers=left,
  numberstyle=\tiny\color{gray},
  showstringspaces=false,
  columns=fullflexible,
  keepspaces=true
}

\newcommand{\perc}{\mathrm{perc}}
\newcommand{\self}{\mathrm{self}}
\newcommand{\file}[1]{\texttt{#1}}

\title{Adding a \texttt{percRecip} Effect to the Three-Way SAOM\\[4pt]
\large Implementation Log --- April 30, 2026}
\author{Jinwoo Cho}
\date{}

\begin{document}
\maketitle

%% ==========================================================
\section{Overview}
%% ==========================================================

This note records every change made to the RSiena fork to add the
\emph{perceived reciprocity} effect (\texttt{percRecip}) and to set up a
parameter-recovery study for it.  Two layers of work were involved:

\begin{enumerate}[leftmargin=*]
  \item A new C++ effect class plus its registration in the build
        wiring.  This is required because the statistic
        $\sum_i \sum_j x_{iij} x_{iji}$ is not expressible as a
        \texttt{crprod} of two existing dependent networks --- the two
        ``diagonals'' it pairs come from the same simulated three-way
        object, so the existing effect framework cannot evaluate the
        contribution function without code-level support.
  \item An R-side recovery harness driven by
        \texttt{siena07(simOnly = TRUE)} that exercises the new effect
        end-to-end against a known truth.  Several iterations were
        needed to align the harness with the fork's three-way mode.
\end{enumerate}

\paragraph{Statistic.}
For perceiver $i$, sender $j$, receiver $k$, define the per-slice
statistic
\[
s_i \;=\; \sum_{k} x_{i,i,k}\, x_{i,k,i},
\]
i.e., the count of alters $m$ such that perceiver $i$ believes there is a
mutual tie between $i$ and $m$.  Aggregated across slices,
\[
S \;=\; \sum_{i} s_i \;=\; \sum_{i}\sum_{k} x_{iik}\, x_{iki},
\]
which is the perceived reciprocity statistic.

%% ==========================================================
\section{C++ Effect Class}
%% ==========================================================

Two new source files implement the effect, modeled on the existing
\file{ThreeWayPerceiverEffect} pair:

\begin{itemize}[leftmargin=*]
  \item \file{src/model/effects/ThreeWayPerceivedReciprocityEffect.h}
  \item \file{src/model/effects/ThreeWayPerceivedReciprocityEffect.cpp}
\end{itemize}

The class derives directly from \texttt{NetworkEffect} (no covariate
dependence), parses the perceiver index $i$ from the variable name
\texttt{Y[n]} at \texttt{initialize()} time, and defines the two
required virtual methods.

\paragraph{Contribution function.}
With \texttt{ego} $= j$, \texttt{alter} $= k$, and perceiver $i$ stored
as \texttt{lperceiverIndex},
\[
\Delta s_i \;=\;
\begin{cases}
P^{(i)}_{k,i} & \text{if } j = i \quad \text{(toggling } i\!\to\!k\text{)}\\[2pt]
P^{(i)}_{i,j} & \text{if } k = i \quad \text{(toggling } j\!\to\!i\text{)}\\[2pt]
0 & \text{otherwise.}
\end{cases}
\]

The corresponding code:

\begin{lstlisting}[language=C++]
double ThreeWayPerceivedReciprocityEffect::calculateContribution(
    int alter) const
{
    const int ego = this->ego();
    const int i   = this->lperceiverIndex;

    if (ego == i) {
        return this->inTieExists(alter) ? 1.0 : 0.0;
    }
    if (alter == i) {
        return this->pNetwork()->tieValue(i, ego) != 0 ? 1.0 : 0.0;
    }
    return 0.0;
}
\end{lstlisting}

\paragraph{Evaluation statistic.}
\texttt{tieStatistic(alter)} sums to $s_i$ when invoked over every
existing tie of the slice.  To avoid double-counting each reciprocated
pair (visited once as $(i,m)$ and once as $(m,i)$), the routine returns
$1$ only on the leg where \texttt{ego} $=$ \texttt{lperceiverIndex}:

\begin{lstlisting}[language=C++]
double ThreeWayPerceivedReciprocityEffect::tieStatistic(int alter)
{
    const int ego = this->ego();
    const int i   = this->lperceiverIndex;
    if (ego == i && this->inTieExists(alter))
        return 1.0;
    return 0.0;
}
\end{lstlisting}

%% ==========================================================
\section{Build Wiring}
%% ==========================================================

The new files must be visible to the compiler, the effect factory, the
package's master include, and the metadata table that drives the R
side.  Five files were touched.

\paragraph{1.\ \file{src/model/effects/AllEffects.h}.}
Added
\verb|#include "ThreeWayPerceivedReciprocityEffect.h"| at the bottom of
the three-way SAOM block, alongside the existing perceiver-attribute
effects.

\paragraph{2.\ \file{src/model/effects/EffectFactory.cpp}.}
A new branch in the \texttt{createEffect()} dispatch maps the short
name to the C++ class:
\begin{lstlisting}[language=C++]
else if (effectName == "percRecip")
{
    pEffect = new ThreeWayPerceivedReciprocityEffect(pEffectInfo);
}
\end{lstlisting}
Inserted directly after the existing \texttt{percSenderAgree} branch.

\paragraph{3.\ Source-list files.}
The package uses a configure-time substitution from
\file{src/sources.list} into \file{src/Makevars.in}.  We appended
\texttt{model/effects/ThreeWayPerceivedReciprocityEffect.cpp} to all
three:
\begin{itemize}[leftmargin=*]
  \item \file{src/sources.list}
  \item \file{src/Makevars}
  \item \file{src/Makevars.win}
\end{itemize}
The first is the canonical source; the second is regenerated by
\file{configure}; the third is the Windows build file.

\paragraph{4.\ \file{data/allEffects.csv}.}
A new metadata row registers the effect for the R side:
\begin{lstlisting}
nonSymmetricObjective,perceived reciprocity,Number of perceiver-involved reciprocated ties,percRecip,TRUE,,,eval,FALSE,FALSE,FALSE,FALSE,FALSE,",",0,0,objective,NA,NA,0,0,0,0,dyadic,TRUE
\end{lstlisting}
Inserted directly after the standard \texttt{recip} row.  The companion
\file{data/allEffects.R} reads the CSV unchanged, so no additional
edits there.

%% ==========================================================
\section{R Recovery Harness}
%% ==========================================================

The recovery script \file{recovery\_percRecip.R} (project root) follows
the standard two-stage pattern:

\begin{enumerate}[leftmargin=*]
  \item One \texttt{siena07(simOnly = TRUE, n3 = 50, returnDeps = TRUE)}
        call draws 50 wave-2 realizations under known truth.
  \item A 50-iteration estimation loop fits each replicate with a
        fresh effects object and stores $\hat\theta$ together with its
        SE.
\end{enumerate}

\subsection{Final pilot configuration}

\begin{itemize}[leftmargin=*]
  \item $n = 10$ actors, hence $K = 10$ perceiver slices.
  \item Two waves, $50$ replications.
  \item \texttt{shareParameters = TRUE} on the threeway dependent
        variable: density, reciprocity, and \texttt{percRecip} share a
        single $\beta$ across slices, materialized at \texttt{Y[shared]}.
  \item \texttt{Y[self]}: density and reciprocity (no cross-network
        \texttt{crprod} effects in the pilot).
  \item Rate parameters set explicitly (see Section~\ref{sec:rates}).
\end{itemize}

Truth values:
\[
\rho_\perc = 6,\quad \rho_\self = 6,\quad
\beta_{\text{dens}}^\perc = -2,\quad
\beta_{\text{recip}}^\perc = 1.5,\quad
\beta_{\text{percRecip}} = 0.8,\quad
\beta_{\text{dens}}^\self = -2,\quad
\beta_{\text{recip}}^\self = 1.5.
\]

\subsection{Iterations during debugging}

Several issues surfaced while the script was being aligned with the
fork's three-way API.  They are recorded here because each one masked
the next.

\paragraph{(a) Misuse of \texttt{sienaDependent(name = \dots)}.}
Initial drafts called
\verb|sienaDependent(arr, name = "Y[1]")|, which errored with
``unused argument''.  \texttt{sienaDependent} has no \texttt{name}
parameter; \texttt{sienaDataCreate} reads variable names from its own
call expression via \texttt{substitute()}.  An interim fix used a
\texttt{do.call} with a named list to bind the non-syntactic name
\texttt{"Y[1]"}.  This was later replaced (see (b)) by the fork's
native three-way constructor.

\paragraph{(b) Wrong data construction for three-way mode.}
The fork's \texttt{sienaDependent} accepts a 4D array of shape
$(K,n,n,T)$ when \texttt{type = "threeway"} is supplied; the slice
names and the self-net are derived automatically.  The harness was
rewritten to use this API:

\begin{lstlisting}[language=R]
Y   <- sienaDependent(arr_4d, type = "threeway",
                      shareParameters = TRUE)
dat <- sienaDataCreate(Y)
eff <- getEffects(dat)
\end{lstlisting}

\paragraph{(c) Canonical row name under \texttt{shareParameters}.}
\texttt{R/effects.r} (line 1007) renames the first occurrence of every
shared structural row from \texttt{Y[1]} to \texttt{Y[shared]}.
Calls of the form
\verb|setEffect(eff, density, name = "Y[1]", ...)|
therefore failed with ``Effect not found''.  The fix replaces
\texttt{Y[1]} with \texttt{Y[shared]} for all canonical-slice
\texttt{setEffect} and \texttt{includeEffects} calls.

\paragraph{(d) Spurious \texttt{fix = TRUE} for simulation.}
\texttt{fix} controls estimation; under \texttt{simOnly = TRUE} no
estimation occurs and \texttt{fix} is irrelevant.  All \texttt{fix =
TRUE} flags were removed; only \texttt{initialValue} is needed to
inject the true parameter values into the simulator.

\paragraph{(e) Rate/density collinearity on \texttt{Y[self]}.}
\label{sec:rates}
The first complete run produced
\begin{quote}\itshape
\textsf{*** Standard errors not reliable *** \\
The following is approximately a linear combination for which the data
carries no information: $1\,\beta_{31} + 1\,\beta_{32}$.}
\end{quote}
The diagnostic \texttt{inc[c(31,32), c("name","shortName")]} identified
the pair as \texttt{Y[self]} rate and \texttt{Y[self]} density.  The
root cause was that no rate parameters had been set in \texttt{truth};
the simulator used a default rate too small to produce identifiable
movement on the self-net.  Two new entries
(\texttt{rate\_perc}, \texttt{rate\_self}) were added to \texttt{truth}
and fed in via a per-slice loop:
\begin{lstlisting}[language=R]
for (k in seq_len(n_actors)) {
  sim_eff <- setEffect(sim_eff, Rate, type = "rate",
                       name = paste0("Y[", k, "]"),
                       period = 1,
                       initialValue = truth["rate_perc"])
}
sim_eff <- setEffect(sim_eff, Rate, type = "rate",
                     name = "Y[self]", period = 1,
                     initialValue = truth["rate_self"])
\end{lstlisting}
The loop is necessary because \texttt{shareParameters} only
share-marks \emph{structural} rows (\texttt{interaction1 == ""}); rate
rows remain per-slice.

\paragraph{(f) Silent name mismatch in the estimation loop.}
\texttt{build\_rep\_data} originally bound the per-replicate dependent
variable to the symbol \texttt{Yrep}, so the resulting slice names
became \texttt{Yrep[shared]}, \texttt{Yrep[self]}, etc.  The loop's
\verb|includeEffects(..., name = "Y[shared]")| therefore matched no
row.  \texttt{includeEffects} does not stop on a missing name --- it
emits a \texttt{cat} message and returns the table unchanged
(\file{R/sienaeffects.r} lines 142--153).  Across 50 replications all
\texttt{percRecip} and \texttt{recip} additions silently no-opped,
yielding 50 \texttt{percRecip-hat = NA} estimates.  The fix is one
line: bind the variable to \texttt{Y} inside \texttt{build\_rep\_data},
matching the base data.

%% ==========================================================
\section{Files Modified or Created}
%% ==========================================================

\begin{tabular}{@{}lll@{}}
\toprule
File & Status & Purpose \\
\midrule
\file{src/model/effects/ThreeWayPerceivedReciprocityEffect.h}
  & new & class declaration \\
\file{src/model/effects/ThreeWayPerceivedReciprocityEffect.cpp}
  & new & class implementation \\
\file{src/model/effects/AllEffects.h}
  & edit & include the new header \\
\file{src/model/effects/EffectFactory.cpp}
  & edit & dispatch \texttt{percRecip} short name \\
\file{src/sources.list}
  & edit & add new \texttt{.cpp} to canonical list \\
\file{src/Makevars}
  & edit & regenerated from \file{Makevars.in}; updated for parity \\
\file{src/Makevars.win}
  & edit & Windows source list \\
\file{data/allEffects.csv}
  & edit & metadata row for \texttt{percRecip} \\
\file{recovery\_percRecip.R}
  & new & two-stage recovery harness \\
\file{notes/percRecip\_changes.tex}
  & new & this note \\
\bottomrule
\end{tabular}

%% ==========================================================
\section{Status}
%% ==========================================================

\begin{itemize}[leftmargin=*]
  \item C++ effect compiles cleanly under syntax-only check
        (\texttt{g++ -fsyntax-only}) against the package include path.
  \item Full installation via \texttt{devtools::load\_all()} succeeds
        on the user's machine.
  \item Simulation runs and produces valid wave-2 networks across all
        50 replications.
  \item Estimation loop now correctly includes \texttt{percRecip} and
        the explicit rate parameters; the rate/density collinearity
        warning has been addressed by setting non-default rates in
        \texttt{truth}.
  \item Pending: a clean recovery summary table showing recovered
        $\hat\beta_{\text{percRecip}} \approx 0.8$ with reasonable
        coverage.  After confirmation, the pilot lifts to a multi-slice
        run with cross-network \texttt{crprod} effects re-enabled.
\end{itemize}

\end{document}
